\section{Final Specifications and Requirements}

The framework is designed to produce three outputs: a body condition estimate, a behavior label, and an identity representation for cattle retrieval. It is evaluated as a research prototype, not as a finished farm product. Because of this, the requirements include both prediction functions and the controls needed to interpret the results properly. These include clear label definitions, protected data splits, failure handling, task-specific metrics, and reproducible settings.

\begin{table}[htbp]
\centering\small
\caption[Research requirements]{Research requirements and the evidence needed for acceptance.}\label{tab:requirements}
\begin{tabularx}{\textwidth}{lYY}
\toprule
ID & Requirement & Evidence required \\
\midrule
R1 & BCS output with explicit units & Physical-score mapping, output support and saved metrics. \\
R2 & Target-cow Behavior output & Taxonomy, crop/cache integrity and class-wise results. \\
R3 & Re-ID under a named protocol & Identity partition, query/gallery definition and retrieval artifact. \\
R4 & Evaluated integrated model & Matched E0/E1/E3 comparisons and sharing boundaries. \\
R5 & Provenance-aware evaluation & Grouped splits, versioned checks and limitations. \\
R6 & Reproducible execution & Executed settings, seed, revision, hashes and checkpoint. \\
R7 & Transparent failures & Missing-input/perception counts; no silent substitutions. \\
R8 & Feasible implementation & Measured resources and maintainable dependencies. \\
\bottomrule
\end{tabularx}
\end{table}


For BCS, the model must predict within the documented score range. The behavior model must use the fixed five-class label set. The Re-ID model must clearly define the learning identities, gallery, queries, and retrieval metrics. Requirement R4 covers the final integration of all three tasks into one multi-task model, including a clear description of what is shared and what remains task-specific.

The main non-functional requirements are reproducibility, leakage resistance, maintainability, clear failure reporting, and practical execution. Dataset manifests, scripts, metrics, and configurations are versioned. Raw datasets and large checkpoints are stored separately. Reproducing an experiment therefore requires both the recorded configuration and access to the correct data.

Scientific validity adds further requirements. Test data must not be used for checkpoint selection. If perception failures remove samples, comparisons must either use matched samples or report the coverage clearly. Missing task labels must be handled with task-specific batches or another available-label method, not treated as negative labels. Finally, every evaluation claim must match the unit protected by the split, such as a burst group, source video, recording session, or biological identity.

\section{Societal Impact}

A unified visual monitoring system could help organize repeated cattle observations by linking body condition and behavior to individual animals. In practice, this could help people notice unusual changes earlier and reduce the amount of manual video checking. However, this study does not include farmer interviews, a field deployment, or measurements of labor savings, productivity, treatment outcomes, or animal welfare improvement. These remain possible benefits rather than measured results.

The model should support human decision-making rather than replace it. A wrong identity match could attach observations to the wrong cow, while an incorrect BCS or behavior prediction could delay attention. A practical system should therefore keep the source image, time, confidence or retrieval evidence, and a way for a person to review the output. Veterinary and welfare decisions still require proper professional judgment.

Accessibility also depends on the farm environment. Camera placement, housing, breed, lighting, network access, and staff practices can vary widely. Results from the selected datasets cannot prove that the system will work in every production setting, including farms in Bangladesh. Testing with local users and local conditions would be needed before making claims about real-world usability or adoption.

Farm images may also contain people, layouts, or commercially sensitive information. Permission to use a dataset for research does not automatically allow redistribution or commercial use. Privacy, ownership, and licensing conditions must therefore be checked for every dataset and for any future field deployment.

\section{Environmental Impact}

The computational cost of the system includes perception preprocessing, model training, data storage, and data transfer. It is not limited to the final prediction network. Even a small downstream model can depend on a larger detector or segmentation model. For this reason, any efficiency claim should make clear whether upstream processing is included.

We reduced repeated computation by checking the pipeline before full experiments, reusing validated cached outputs, and using deterministic manifests. These steps reduce unnecessary repeated processing, but we did not collect a complete power trace or location-specific emissions measurement. Therefore, the study does not claim measured carbon savings or measured emissions.

Better monitoring could eventually support more efficient livestock management. However, the experiments do not measure feed savings, methane reduction, or changes in disease burden. Those effects would depend on how the predictions change real farm decisions and would need a separate long-term study.

\section{Ethical Issues}

A major ethical responsibility in this work is to separate what the data actually show from what we might assume. ScienceDB is not described as cow-disjoint because verified biological cow IDs are not available. The primary behavior evaluation is not described as unseen-cow evaluation. The oracle segmentation-guided Re-ID experiment uses ground-truth masks, so it is not evidence of an automatic segmentation system. Positive, mixed, and negative findings are all reported.

Dataset reuse must also be checked for each source. The dataset registry records the available licenses and access conditions, but any unresolved permission must be confirmed before redistribution or deployment. The thesis does not claim institutional ethics approval, an exemption, a user study, or field consent unless a supporting institutional record exists.

Validation data are used for model selection, while held-out test data are reserved for evaluation. External test data must not be used for tuning after the results are seen. Samples removed by perception failures are reported through coverage information instead of being silently ignored.

The authors are still responsible for work produced with AI-assisted writing or coding. They must review the final manuscript, verify the claims, follow university disclosure rules, and confirm the declaration and acknowledgments.

% TODO-P3-ADMIN: Confirm permissions and institutional requirements.
\drafttodo{ETHICS-02}{Confirm dataset-specific reuse permissions and the university's ethics and AI-assistance disclosure requirements. No approval or exemption is asserted in advance.}

\section{Standards and Technical Conventions}

The study follows documented technical conventions, but it does not claim certification under an engineering standard that has not been verified. BCS labels are interpreted using their released score range and established anatomical meaning~\cite{edmonson1989bcs,ferguson1994bcs}. Behavior labels use a versioned mapping, and Re-ID uses a defined query--gallery protocol. These rules make the evaluation more consistent, but they do not mean that the system is formally certified.

The thesis follows the supplied BRAC University chapter structure and uses an IEEE-style bibliography. Dataset licenses, privacy requirements, deployment safety, and any separate IEEE-format submission requirement must be checked independently. Using IEEE reference formatting does not mean that the work has been submitted to or accepted by a journal.

\section{Project Management Plan}

The project is organized into several main stages: protocol verification, perception assessment, single-task reference models, cattle-centered alternatives, and multi-task comparison. Each stage provides the information needed for the next. Table~\ref{tab:timeline} records the main project milestones. The timeline shows project progress, not scientific evidence.

\begin{table}[htbp]
\centering\small
\caption[Project milestones and remaining work]{Major Phase 3 milestones and the remaining MTL boundary.\evidence{E01,E02,E03,E23,E25,E28,E30}}\label{tab:timeline}
\begin{tabularx}{\textwidth}{lYY}
\toprule
Stage & Milestone & Status \\
\midrule
Earlier phases & Problem formulation and prior thesis work & Historical context \\
20 Sep 2026 & ScienceDB repair and SideView protocols & Recorded \\
20--22 Sep & Perception and visual audits & Partial feasibility \\
23 Sep & CVB plus Beef protocol and RGB results & Recorded \\
23 Sep & RGB baselines and SideView hydration & Runs 1--3 completed \\
24 Sep & Cattle-centered task models & Runs 4--6 completed and evaluated \\
Next experiment & E1 hard-shared MTL control & Pending Run 7 \\
Final experiment & E3 modular/task-private MTL & Pending Run 8 \\
After Runs 7--8 & Integrated discussion, abstract, conclusion & Pending final synthesis \\
\bottomrule
\end{tabularx}
\end{table}


Computational resources were selected according to the needs of preprocessing, training, and evaluation. Small pipeline checks were done before full experiments, and reusable intermediate outputs reduced repeated preprocessing. Detailed hardware and execution records are kept for reproducibility instead of being included as part of the main scientific argument.

Versioned code, deterministic manifests, persistent data storage, and machine-readable results connect each experiment to its settings. These records help keep execution consistent across different environments and provide the information needed to reproduce reported results.

Author contributions, author order, committee information, and submission details still need confirmation from the authors and supervisors. They are not guessed from repository activity.

\section{Risk Management}

Table~\ref{tab:risks} lists the main scientific and operational risks. These include related observations crossing train and test splits, unsupported identity claims, classes linked to one data source, perception failures, incorrect causal claims from combined model changes, and unfair comparisons caused by different training budgets.

\begin{table}[htbp]
\centering\small
\caption[Research risks and mitigations]{Scientific and operational risks; implemented controls do not eliminate all residual risks.\evidence{E04,E05,E06,E07,E11,E24,E26,E29}}\label{tab:risks}
\begin{tabularx}{\textwidth}{lYY}
\toprule
Risk & Consequence & Control / open issue \\
\midrule
Correlated data & Optimistic evaluation & Burst/source/session grouping; identity may remain unknown. \\
Temporary IDs & False cross-cow claim & Use biological IDs only where verified. \\
Unavailable data & Blocked experiments & Approved SideView contingency; MultiCam still blocked. \\
Source/class link & Shortcut prediction & Source/class metrics; Walking is CVB-only. \\
Perception error & Missing/wrong target input & Record coverage and exclude invalid masks without fabrication. \\
Combined changes & Unsupported causal claim & Attribute results to the full executed configuration. \\
Oracle masks & Overstated deployment readiness & Label Run 6 as GT-guided, not automatic segmentation. \\
Single-run results & Unknown training variability & Report point estimates; do not invent significance. \\
External mismatch & Misleading transfer & Freeze defensible mappings before testing. \\
Limited resources & Incomplete comparisons & Local smoke checks, focused cloud training, and explicit deferral of unexecuted E2/E5 and external stress tests. \\
Unclear license & Uncertain reuse rights & Confirm permission; do not infer from access. \\
\bottomrule
\end{tabularx}
\end{table}


The evaluation design reduces these risks using task-specific controls. ScienceDB burst groups and behavior videos or sessions stay inside their assigned partitions. SideView Protocol A separates biological identities between representation learning and evaluation. Matched subsets are used to compare the RGB and perception-enhanced BCS and behavior models on the same retained samples. The Re-ID model is clearly described as oracle-guided when it uses ground-truth masks. Important limits still remain: ScienceDB cow identities are unknown, Walking is linked to one source, and the task models are based on single runs. The final MTL comparison must keep the task protocols comparable and report any difference in training budget.

\section{Economic Analysis}

The available records allow only a partial discussion of research cost, not a full farm-deployment business case. The experiments use different accelerators, preprocessing steps, storage needs, and billing conditions. Their runtime therefore cannot be turned into directly comparable cost values without consistent rates or invoices. Table~\ref{tab:costs} reports only information that was actually recorded.

\begin{table}[htbp]
\centering\small
\caption[Recorded runtime and partial costs]{Selected recorded training/pipeline runtimes; these are not a total project or deployment budget.\evidence{E08,E09,E14,E22,E24,E28,E29}}\label{tab:costs}
\begin{tabularx}{\textwidth}{lYY}
\toprule
Item & Elapsed time & Cost evidence \\
\midrule
BCS baseline & 2,023.22 s & About USD 1.15 in research log; not invoice-reconciled. \\
Behavior baseline & 821.47 s & No dollar total in inspected result JSON. \\
Re-ID RGB baseline & 1,049.05 s training & No dollar total recorded in the artifact. \\
BCS cattle-centered & 550.80 s across recorded epochs & Perception preprocessing billed separately; no total here. \\
Behavior cattle-centered & 1,009.91 s training & Cache generation was a separate 3,492.3 s job. \\
Re-ID oracle GT-guided & 2,350.67 s full pipeline & Includes training and Protocol A evaluation. \\
Complete system & Not measured & No deployment ROI or energy estimate. \\
\bottomrule
\end{tabularx}
\end{table}


A real deployment estimate would need information such as the number and placement of cameras, how often inference runs, the required computing hardware, network access, storage period, maintenance, and human-review workload. These values are not fixed by the benchmark experiments. Therefore, the study does not claim a measured return on investment, a payback period, or proven cost savings from multi-task integration.
